% ===== PRÉAMBULE COMMUN — À COPIER TEL QUEL EN TÊTE DE CHAQUE FICHIER .tex =====
% Ne modifier QUE :
%   - les deux lignes \fancyhead[L] et \fancyhead[R]
%   - le bloc titre \begin{center}...\end{center} juste après \begin{document}
\documentclass[11pt,a4paper]{article}
\usepackage[utf8]{inputenc}
\usepackage[T1]{fontenc}
\usepackage{lmodern}
\usepackage[french]{babel}
\usepackage{amsmath,amssymb,mathtools}
\usepackage{icomma}
\usepackage{xcolor}
\usepackage{tikz}
\usepackage{pgfplots}
\usepackage[most]{tcolorbox}
\usepackage{enumitem}
\usepackage{array,booktabs,colortbl}
\usepackage[a4paper,margin=2.2cm,headheight=26pt,headsep=10pt]{geometry}
\usepackage{fancyhdr}
\usepackage[hidelinks]{hyperref}
\usetikzlibrary{arrows.meta,calc,angles,quotes,positioning,decorations.markings,intersections,backgrounds}
\pgfplotsset{compat=1.18}

\definecolor{primary}{RGB}{226,74,88}
\definecolor{primarydark}{RGB}{150,28,42}
\definecolor{methcol}{RGB}{40,90,150}
\definecolor{excol}{RGB}{0,135,90}
\definecolor{defcol}{RGB}{0,114,178}
\definecolor{attncol}{RGB}{200,40,55}
\definecolor{thmcol}{RGB}{0,140,120}
\definecolor{courbe}{RGB}{32,110,200}
\definecolor{courbeder}{RGB}{210,70,40}
\definecolor{tang}{RGB}{0,150,90}
\definecolor{grille}{RGB}{200,205,215}
\definecolor{plus}{RGB}{200,60,60}
\definecolor{moins}{RGB}{30,90,200}

\tcbset{boxbase/.style={enhanced,breakable,boxrule=0.9pt,arc=2.5pt,
  left=3mm,right=3mm,top=2.3mm,bottom=2.3mm,fonttitle=\bfseries,coltitle=white}}
\newtcolorbox{cle}[1][]{boxbase,colback=defcol!6,colframe=defcol,
  title={\textsc{À retenir}\ifx\relax#1\relax\relax\else\textmd{ : \itshape #1}\fi}}
\newtcolorbox{methode}[1][]{boxbase,colback=methcol!7,colframe=methcol,sharp corners,
  title={\textsc{Méthode}\ifx\relax#1\relax\relax\else\textmd{ --- \itshape #1}\fi}}
\newtcolorbox{exemplebox}[1][]{boxbase,colback=excol!5,colframe=excol,
  title={\textsc{Exemple}\ifx\relax#1\relax\relax\else\textmd{ : \itshape #1}\fi}}
\newtcolorbox{piege}[1][]{boxbase,colback=attncol!6,colframe=attncol,
  title={\textsc{Attention}\ifx\relax#1\relax\relax\else\textmd{ : \itshape #1}\fi}}
\newtcolorbox{exo}[1][]{boxbase,colback=primary!4,colframe=primary,
  title={\textsc{Exercice}\ifx\relax#1\relax\relax\else\textmd{~#1}\fi}}
\newtcolorbox{corr}[1][]{boxbase,colback=thmcol!5,colframe=thmcol,
  title={\textsc{Corrigé}\ifx\relax#1\relax\relax\else\textmd{~#1}\fi}}

\newcommand{\R}{\mathbb{R}}
\newcommand{\N}{\mathbb{N}}
\newcommand{\ds}{\displaystyle}
\newcommand{\tg}{\operatorname{tg}}
\newcommand{\cotg}{\operatorname{cotg}}
\newcommand{\arctg}{\operatorname{arctg}}
\newcommand{\arccotg}{\operatorname{arccotg}}
\newcommand{\limh}{\lim_{h\to 0}}

\pagestyle{fancy}\fancyhf{}
\fancyhead[L]{\small\textcolor{primarydark}{\textbf{Thème 1} — Nombre dérivé \& taux d'accroissement}}
\fancyhead[R]{\small\textcolor{primarydark}{\textsc{Tutoriel}}}
\fancyfoot[C]{\small\thepage}
\renewcommand{\headrulewidth}{0.8pt}

\begin{document}

\begin{center}
{\LARGE\bfseries\color{primarydark} Nombre dérivé \& taux d'accroissement}\\[2pt]
{\color{primary}\rule{0.5\textwidth}{1.2pt}}\\[4pt]
{\small\itshape De la pente de la sécante à la pente de la tangente : tout par la définition}
\end{center}
\vspace{2mm}

\section*{1. Le taux d'accroissement : une variation moyenne}

Soit $f$ une fonction et $x_0$ un point de son domaine. Quand on part de $x_0$ et qu'on
avance d'un petit pas $h$ (avec $h\neq 0$), la fonction passe de $f(x_0)$ à $f(x_0+h)$.
Le \textbf{taux d'accroissement} de $f$ entre $x_0$ et $x_0+h$ mesure la variation
\emph{moyenne} de $f$ sur cet intervalle :
\[
\tau(h)=\frac{f(x_0+h)-f(x_0)}{h}
=\frac{\text{variation de la fonction}}{\text{variation de la variable}}.
\]

\begin{cle}[trois lectures d'un même nombre]
Le taux d'accroissement $\tau(h)$ se lit de trois façons équivalentes :
\begin{itemize}[leftmargin=6mm,itemsep=1pt]
  \item \textbf{Variation moyenne} de $f$ entre $x_0$ et $x_0+h$ ;
  \item \textbf{Pente de la sécante} passant par les points $A\big(x_0;f(x_0)\big)$ et
        $B\big(x_0+h;f(x_0+h)\big)$ ;
  \item \textbf{Vitesse moyenne} si $f(x)$ représente une position au temps $x$.
\end{itemize}
\end{cle}

\section*{2. Passage à la limite : le nombre dérivé}

Faisons maintenant tendre le pas $h$ vers $0$ : le point $B$ glisse le long de la courbe
et se rapproche de $A$. La sécante $(AB)$ pivote et tend vers une position limite : la
\textbf{tangente} à la courbe au point $A$. La pente de cette tangente est le
\textbf{nombre dérivé} de $f$ en $x_0$ :
\[
\boxed{\,f'(x_0)=\limh \tau(h)=\limh \frac{f(x_0+h)-f(x_0)}{h}\,}
\]

\begin{center}
\begin{tikzpicture}[scale=1.0,>=Latex]
  \def\a{1.1}   % x0
  \def\b{2.7}   % x0+h
  % axes
  \draw[->,gray!70] (-0.3,0)--(4.2,0) node[right]{\small $x$};
  \draw[->,gray!70] (0,-0.3)--(0,3.7) node[above]{\small $y$};
  % courbe f(x)=0.45 x^2
  \draw[courbe,very thick,domain=-0.2:3.05,smooth,samples=80]
        plot (\x,{0.45*\x*\x});
  \node[courbe] at (3.15,2.9) {\small $\mathcal{C}_f$};
  % points A et B
  \coordinate (A) at (\a,{0.45*\a*\a});
  \coordinate (B) at (\b,{0.45*\b*\b});
  % secante (AB) prolongee
  \draw[primary,thick,shorten >=-1.3cm,shorten <=-1.3cm] (A)--(B);
  % tangente en A : pente 2*0.45*a = 0.9*a
  \draw[tang,very thick,shorten >=-1.4cm,shorten <=-1.2cm]
        (A)++(-1,{-0.9*\a}) -- ++(3,{0.9*\a*3});
  % triangle des accroissements
  \draw[dashed,gray] (A)--(\b,{0.45*\a*\a}) node[midway,below,font=\scriptsize]{$h$};
  \draw[dashed,gray] (\b,{0.45*\a*\a})--(B)
        node[midway,right,font=\scriptsize]{$f(x_0{+}h){-}f(x_0)$};
  % points
  \fill[primarydark] (A) circle (1.6pt) node[above left,font=\scriptsize]{$A$};
  \fill[primarydark] (B) circle (1.6pt) node[above left,font=\scriptsize]{$B$};
  % abscisses
  \draw[gray!70] (\a,0.05)--(\a,-0.05) node[below,font=\scriptsize]{$x_0$};
  \draw[gray!70] (\b,0.05)--(\b,-0.05) node[below,font=\scriptsize]{$x_0{+}h$};
  % legendes
  \node[primary,font=\scriptsize] at (3.55,1.15) {sécante};
  \node[tang,font=\scriptsize] at (2.75,0.15) {tangente};
\end{tikzpicture}
\end{center}

\noindent Quand $h\to 0$ : la sécante (rouge) bascule vers la tangente (verte), et la
variation moyenne $\tau(h)$ devient la \textbf{variation instantanée} $f'(x_0)$.

\begin{cle}[du moyen à l'instantané]
$f'(x_0)$ est à la fois la \textbf{pente de la tangente} en $A$, la \textbf{variation
instantanée} de $f$ en $x_0$, et la \textbf{vitesse instantanée} à l'instant $x_0$ si
$f$ est une position.
\end{cle}

\begin{methode}[calculer $f'(x_0)$ par la définition, en 4 étapes]
\begin{enumerate}[leftmargin=7mm,itemsep=1pt]
  \item \textbf{Écrire} le taux d'accroissement $\ds\tau(h)=\frac{f(x_0+h)-f(x_0)}{h}$.
  \item \textbf{Développer et simplifier le numérateur} $f(x_0+h)-f(x_0)$ (le terme
        constant $f(x_0)$ doit disparaître, laissant $h$ en facteur).
  \item \textbf{Simplifier par $h$} (on a le droit car $h\neq 0$).
  \item \textbf{Faire tendre $h\to 0$} : la limite obtenue est $f'(x_0)$.
\end{enumerate}
\end{methode}

\begin{exemplebox}[la fonction carré $f(x)=x^2$]
\textbf{Étape 1.} $\ds \tau(h)=\frac{f(x_0+h)-f(x_0)}{h}=\frac{(x_0+h)^2-x_0^2}{h}$.\\[2pt]
\textbf{Étape 2.} On développe le numérateur :
\[
(x_0+h)^2-x_0^2=x_0^2+2x_0h+h^2-x_0^2=2x_0h+h^2=h(2x_0+h).
\]
\textbf{Étape 3.} On simplifie par $h$ (licite car $h\neq0$) :
$\ds \tau(h)=\frac{h(2x_0+h)}{h}=2x_0+h$.\\[2pt]
\textbf{Étape 4.} On passe à la limite :
\[
f'(x_0)=\limh (2x_0+h)=2x_0.
\]
Ainsi la pente de la tangente à la parabole au point d'abscisse $x_0$ vaut $2x_0$.
\end{exemplebox}

\begin{piege}
On ne remplace \emph{jamais} $h$ par $0$ dès l'étape 1 : on obtiendrait
$\frac{0}{0}$, forme indéterminée. Il faut d'abord \textbf{simplifier par $h$}
au numérateur, \emph{puis} seulement faire $h\to 0$.
\end{piege}

\end{document}
